% !PRINTED 685
\weberpage{685}
%<<CONT
we can form the linear equations, in analogy with \S.\,186, (5):
\weq{2}{
\begin{aligned}
\xi_1 l_{1,1} + \xi_2 l_{2,1} + \cdots + \xi_{\nu-1} l_{\nu-1,1} + \delta_1 \xi_\nu &= \lambda_1, \\
\xi_1 l_{1,2} + \xi_2 l_{2,2} + \cdots + \xi_{\nu-1} l_{\nu-1,2} + \delta_2 \xi_\nu &= \lambda_2, \\
\cdot\ \cdot\ \cdot\ \cdot\ \cdot\ \cdot\ \cdot\ \cdot\ \cdot\ \cdot\ \cdot\ \cdot & \\
\xi_1 l_{1,\nu} + \xi_2 l_{2,\nu} + \cdots + \xi_{\nu-1} l_{\nu-1,\nu} + \delta_\nu \xi_\nu &= \lambda_\nu,
\end{aligned}
}
whose determinant [\S.\,186, (4)] has the value, different from zero:
\weq{3}{
\left|
\begin{array}{ccccc}
l_{1,1}, & l_{2,1}, & \ldots, & l_{\nu-1,1}, & \delta_1\\
l_{1,2}, & l_{2,2}, & \ldots, & l_{\nu-1,2}, & \delta_2\\
\cdot & \cdot & \cdot & \cdot & \cdot\\
l_{1,\nu}, & l_{2,\nu}, & \ldots, & l_{\nu-1,\nu}, & \delta_\nu
\end{array}
\right| = n\,L\,(\varepsilon_1, \varepsilon_2, \ldots, \varepsilon_{\nu-1})
}
Accordingly the unknowns $\xi_1, \xi_2, \ldots, \xi_{\nu-1}, \xi_\nu$ are
\gesperrt{uniquely} determined from (2).

By addition of the equations (2) one finds first of all,
very simply,
\weq{4}{
n\,\xi_\nu = \log N_a(\alpha),
}
and $\xi_\nu$ therefore remains the same for all \gesperrt{numbers $\alpha$ that are
 associated with one another.}

The numbers $\xi_1, \xi_2, \ldots, \xi_{\nu-1}$ are called the exponents of the
number $\alpha$ [with respect to the fundamental system $(\varepsilon_1, \varepsilon_2, \ldots, \varepsilon_{\nu-1})$].

From this definition it follows at once that the
exponents of a product or of a quotient are equal to the
sum or the difference of the corresponding exponents of the
individual parts. The exponents of a unit are
integral rational numbers, and every system of integral rational numbers
$\xi_1, \xi_2, \ldots, \xi_{\nu-1}$ is also the system of exponents of a unit.

The exponents of associated numbers therefore differ
from one another by integers, and to every number $\alpha$
one can find an associated number $\alpha_0$ whose exponents lie between 0
and 1. Such numbers are called \gesperrt{reduced numbers} (with
respect to the fundamental system $\varepsilon_1, \ldots, \varepsilon_{\nu-1}$). Reduced units
are all and only the \gesperrt{roots of unity} present in $\Omega$,
\gesperrt{whose number we shall denote by $w$.} Since
the two roots of unity $\pm 1$ are contained in every field,
$w$ is at least $= 2$ and always an even number.

If two associated numbers have the same system of exponents,
then they differ only by a factor which
%CONT>>
